\documentclass[11pt,a4paper]{article}
\usepackage[margin=1in]{geometry}
\usepackage{hyperref}
\usepackage{parskip}
\usepackage{mathptmx}

\hypersetup{
    colorlinks=true,
    linkcolor=blue,
    urlcolor=blue,
    citecolor=blue
}

\begin{document}

\begin{raggedleft}
\textbf{Date:} August 5, 2026\\
\end{raggedleft}

\textbf{To:}\\
The Editor-in-Chief\\
\textit{Discover AI}\\
Springer Nature\\

\vspace{1em}

\textbf{Subject: Submission of original research manuscript ``Whose Art Counts? Model- and Prompt-Dependent Associations in Vision-Language Judgments of Museum Art''}

\vspace{1em}

Dear Editor-in-Chief and Editorial Board,

We are pleased to submit our original research manuscript titled \textbf{``Whose Art Counts? Model- and Prompt-Dependent Associations in Vision-Language Judgments of Museum Art''} for consideration for publication as a research article in \textit{Discover AI}.

Vision-language models (VLMs) such as CLIP, OpenCLIP, and SigLIP are increasingly integrated into digital humanities workflows, cultural collection search systems, and automated image retrieval. However, evaluating these models on museum archives presents a core measurement challenge: score disparities can reflect learned image-text associations, visual features, or the underlying structure and selection histories of the catalog itself. 

In this manuscript, we introduce an \textbf{archive-conditioned audit protocol} that evaluates model, prompt, and metadata provenance within a unified measurement framework. We test our protocol on 445 publicly accessible object images from the Metropolitan Museum of Art Open Access collection across four vision-language model architectures and four prompt pairs. Our key findings include:

\begin{enumerate}
    \item \textbf{Model- and Prompt-Dependence}: Prompts framing ``influence'' produce positive male-minus-female score differences for OpenAI CLIP and OpenCLIP ViT-L/14, whereas masterpiece and quality prompts yield inconsistent or non-reproducible patterns across models.
    \item \textbf{Scale-Aware Diagnostics}: Score distributions across model families (e.g., CLIP probability contrasts vs.\ SigLIP logit contrasts) require within-model standardization rather than raw score pooling.
    \item \textbf{Archival Transparency}: By explicitly retaining unknown and unattributed records ($n=384$ of $445$) rather than dropping them, the protocol prevents model evaluations from silently converting into evaluations of named artists only.
\end{enumerate}

We believe \textit{Discover AI} is the ideal venue for this work given the journal's focus on artificial intelligence applications, multimodal representation learning, evaluation methodologies, and responsible AI governance.

\textbf{Declarations \& Compliance}:
\begin{itemize}
    \item This manuscript is original, has not been published previously, and is not currently under consideration for publication elsewhere.
    \item All listed authors have read and approved the submitted version of the manuscript.
    \item The authors declare no competing financial or non-financial interests.
    \item The project repository containing object manifests, prompts, score tables, and complete statistical code is publicly available at \url{https://github.com/manpreet28111995/disentangling-archival-bias}.
\end{itemize}

Thank you very much for considering our manuscript. We look forward to your response.

\vspace{1.5em}

Sincerely,

\vspace{1em}

\textbf{Manpreet Singh} (Submitting Author Author)\\
Boston University, Boston, MA, USA\\
Email: \href{mailto:manni@bu.edu}{manni@bu.edu}\\

\textbf{Rahul Joshi} (Corresponding Author)\\
Symbiosis Institute of Technology, Symbiosis International University, Pune, India\\
Email: \href{rahulj@sitpune.edu.in}{rahulj@sitpune.edu.in}\\



\vspace{0.5em}

\textit{On behalf of co-authors:}\\
\textbf{Nandakishor Reddy Pulagam} (Independent Researcher)\\
\textbf{Rhythm Bhatia} (University of Eastern Finland)\\


\end{document}
